% =====================================================================
% CHAPTER 5: PROPOSED METHODOLOGY
% =====================================================================
\chapter{Proposed Methodology}

\section{Software Development}

\subsection{Control Code Development}
The software architecture will follow a client-server model with the ESP32-CAM acting as the client and the Flask application serving as the backend server. The development process will involve two parallel tracks that will be integrated during the system integration phase.

\subsubsection*{ESP32-CAM Firmware (C++/Arduino)}
The firmware will be developed using the Arduino IDE 2.x with the ESP32 board support package. The firmware will initialize the camera module in JPEG mode at VGA (640$\times$480) resolution, establish Wi-Fi connectivity with the local network using stored SSID and password credentials, and enter a main loop that monitors button inputs via GPIO interrupt handlers.

When the user presses the action button (Button 2), the firmware will execute the following sequence based on the currently selected mode:

\begin{enumerate}[leftmargin=1.5cm]
\item \textbf{Image Capture Mode:} The firmware will capture a single image frame from the OV2640 camera buffer, encode the JPEG data in Base64 format using an embedded Base64 encoding library, construct an HTTP POST request with a JSON payload containing the encoded image and the request type identifier, and transmit it to the Flask server's corresponding endpoint. The total payload size will be approximately 40--70 KB after Base64 encoding.

\item \textbf{Audio Capture Mode:} For speech-to-text functionality, the firmware will sample audio from the INMP441 microphone via the I2S peripheral at 16 kHz sample rate with 16-bit resolution. Audio data will be buffered for a configurable duration (default 5 seconds), then Base64-encoded and transmitted to the \texttt{/transcribe} endpoint.
\end{enumerate}

Upon receiving the JSON response from the server, the firmware will extract the AI-generated text field and render it on the SSD1306 OLED display using the Adafruit SSD1306 and Adafruit GFX libraries. For text responses exceeding the 8-line display capacity, automatic vertical scrolling will be implemented. The response text will also be available for text-to-speech output through the micro speaker.

The mode selection button (Button 1) will cycle through the five operating modes in sequence: Scene Description $\rightarrow$ OCR $\rightarrow$ Voice $\rightarrow$ Q\&A $\rightarrow$ Translate. The current mode will be displayed as a status indicator on the OLED screen.

\subsubsection*{Flask Backend Server (Python)}
The backend will implement five RESTful API endpoints corresponding to the five AI capabilities:

\begin{itemize}[leftmargin=1.5cm]
\item \texttt{POST /analyze-image} --- Scene description using Llama 4 Scout 17B via Groq API
\item \texttt{POST /extract-text} --- Text extraction using local EasyOCR
\item \texttt{POST /transcribe} --- Speech-to-text using local Whisper-tiny
\item \texttt{POST /ask} --- Question answering using Llama 3.3 70B via Groq API
\item \texttt{POST /translate} --- Language translation using Llama 3.3 70B via Groq API
\item \texttt{GET /health} --- Server health check and status endpoint
\end{itemize}

Each endpoint will receive Base64-encoded data in a JSON payload, decode it into the original binary format (JPEG image or WAV audio), and route the request to the appropriate AI model through the Smart Router. The response will be a JSON object containing the AI-generated text, processing time in milliseconds, and a status code.

\subsection{Smart Routing Architecture}
The Smart Router is the core architectural component that will manage model lifecycle and resource allocation:

\begin{itemize}[leftmargin=1.5cm]
\item \textbf{For cloud-based tasks} (scene description, Q\&A, translation): The Smart Router will call the Groq API directly using the Groq Python SDK. Since these models run on Groq's cloud infrastructure, they will consume zero local GPU memory and zero local RAM beyond the HTTP client overhead. The API calls will be asynchronous-capable, with configurable timeout settings and automatic retry logic for transient network failures.

\item \textbf{For local tasks} (OCR, speech-to-text): The Smart Router will first check if the required model is already loaded in memory by querying an internal model registry. If the model is not loaded, the router will first unload any currently loaded local model by deleting the model reference, invoking Python's garbage collector (\texttt{gc.collect()}), and clearing the PyTorch CUDA memory cache (\texttt{torch.cuda.empty\_cache()}). Only then will the required model be loaded into GPU memory. After processing the request, the model will be explicitly unloaded using the same cleanup procedure to free GPU memory for subsequent requests.
\end{itemize}

This \textit{load-on-demand, unload-after-use} strategy will ensure that the system never has more than one local AI model resident in GPU memory at any given time, keeping peak memory usage well below the capacity of entry-level NVIDIA GPUs.

\section{Data Flow Pipeline}
The data flow through the system will follow a four-stage pipeline, illustrated as follows:

\begin{enumerate}[leftmargin=1.5cm]
\item \textbf{Stage 1 --- Capture:} The ESP32-CAM will capture a JPEG image at 640$\times$480 resolution (approximately 30--50 KB after compression) using the OV2640 camera module, or record audio from the INMP441 microphone at 16 kHz sample rate with 16-bit resolution for a configurable duration. The capture will be triggered by the user pressing Button 2.

\item \textbf{Stage 2 --- Encode and Transmit:} The captured data will be Base64-encoded to ensure safe transmission over HTTP without data corruption. The encoded data will be wrapped in a JSON payload with metadata fields indicating the request type, operating mode, and timestamp. The JSON payload will be transmitted as an HTTP POST request body to the Flask server's corresponding endpoint over the local Wi-Fi network.

\item \textbf{Stage 3 --- Process (Smart Routing):} The Flask Smart Router will receive the HTTP request, validate the JSON payload format, decode the Base64 data, and dispatch it to the appropriate AI model based on the endpoint called. Cloud API calls (Groq) are expected to complete in 0.4--0.7 seconds, while local model processing (EasyOCR, Whisper) is expected to take 0.4--1.8 seconds depending on the complexity of the input data.

\item \textbf{Stage 4 --- Respond and Display:} The AI-generated text response will be formatted as a standardized JSON object containing the response text, processing time, and status code. The JSON response will be sent back to the ESP32-CAM via HTTP. The ESP32 firmware will parse the JSON response, extract the text field, render it on the OLED display with word wrapping, and optionally convert it to speech for audio output through the micro speaker.
\end{enumerate}

\section{Testing Plan}
The testing strategy will employ a three-tier approach to systematically validate all aspects of the system before deployment:

\subsection{Tier 1: AI Model Tests}
Each AI model will be tested independently to verify correct functionality. The following tests will be conducted:

\begin{itemize}[leftmargin=1.5cm]
\item \textbf{Llama 4 Scout 17B (Vision):} Will be tested with sample images containing diverse scenes (indoor, outdoor, text-heavy, object-rich) to verify that the model generates accurate and detailed scene descriptions via the Groq API.
\item \textbf{Llama 3.3 70B (Q\&A):} Will be tested with factual questions across multiple domains to verify accurate and contextually relevant answers.
\item \textbf{Llama 3.3 70B (Translation):} Will be tested with English-to-Hindi and English-to-Kannada translation tasks to verify linguistic accuracy.
\item \textbf{EasyOCR:} Will be tested with images containing printed text in various fonts, sizes, and orientations to verify text extraction accuracy.
\item \textbf{Whisper-tiny:} Will be tested with audio samples of varying duration, accent, and background noise levels to verify transcription accuracy.
\end{itemize}

\subsection{Tier 2: Memory Management Tests}
Memory management tests will verify the Smart Routing architecture's resource optimization:

\begin{itemize}[leftmargin=1.5cm]
\item Groq API calls will be verified to consume zero local GPU memory.
\item Local models (EasyOCR, Whisper) will be verified to load and unload cleanly without memory leaks across multiple load-unload cycles.
\item The Smart Router will be tested under mixed workloads to verify that peak GPU memory usage remains below the 600 MB target threshold.
\end{itemize}

\subsection{Tier 3: Flask Endpoint Integration Tests}
Each API endpoint will be tested with valid input data and validated for correct HTTP response codes (200 for success, 400 for malformed requests, 500 for server errors) and properly formatted JSON responses. Error handling will be verified by submitting malformed requests (invalid Base64 data, missing fields, oversized payloads) and checking for appropriate error responses with descriptive error messages.
